% ============================================================
%  HCMUT-DEE Thesis Kit v2.0.4 — Abstract (English)
%  Copyright (c) 2026 Nguyễn Trọng Thắng & TS. Nguyễn Phúc Khải
%  License: MIT
% ============================================================

\chapter*{ABSTRACT}
\addcontentsline{toc}{chapter}{Abstract}

Drafting technical reports and graduation theses presents a significant challenge for engineering students due to strict academic formatting requirements. Traditional word processors, such as Microsoft Word, often exhibit severe limitations regarding document consistency, formatting glitches across different platforms, and difficulties in managing source code, figures, and bibliographies. To address these issues, this project establishes the **HCMUT-DEE Thesis Kit**, an optimized LaTeX template system specifically engineered for students of the Faculty of Electrical and Electronic Engineering (FEEE) at Ho Chi Minh City University of Technology (HCMUT - VNU-HCM).

This document serves as both a comprehensive user guide and the official structural sample of the template. The toolkit is developed around a robust document class (`hcmut-dee.cls`) derived from the standard `report` class, integrating essential packages for advanced mathematical formatting, automated cover page (\texttt{\textbackslash makecover}), official assignment sheet (\texttt{\textbackslash makeassignment}), and grading form (\texttt{\textbackslash makeevaluation}) generation. Furthermore, the template provides pre-configured code listing styles for Python, MATLAB, C/C++, and PLC Structured Text, fully satisfying the specialized academic needs of Electrical Power, Electronics, Telecommunications, and Control Engineering majors.

Practical applications demonstrate that the template reduces typesetting time by over 80\%, completely eliminates overflow margins and citation format errors, and significantly elevates the typographic and professional standards of academic reports. This project embodies the spirit of knowledge heritage among student generations and aims to promote the wide adoption of LaTeX in engineering and scientific research at the university.

\vspace{0.5cm}
\noindent \textbf{Keywords:} LaTeX Template, Technical Report, Graduation Thesis, Faculty of Electrical and Electronic Engineering, HCMUT, Knowledge Heritage.
